\section{Marco teórico}

Un péndulo simple se modela como una masa puntual suspendida de un hilo inextensible y de masa despreciable, de longitud $L$. Cuando la masa se separa de su posición de equilibrio un ángulo $\theta$, la componente tangencial del peso actúa como fuerza restauradora y se expresa como

\begin{equation}
F_t=-mg\sin\theta.
\label{eq:fuerza_restauradora}
\end{equation}

Aplicando la segunda ley de Newton al movimiento tangencial del sistema se obtiene

\begin{equation}
mL\ddot{\theta}=-mg\sin\theta,
\label{eq:segunda_ley_pendulo}
\end{equation}

de donde resulta la ecuación de movimiento

\begin{equation}
\ddot{\theta}+\frac{g}{L}\sin\theta=0.
\label{eq:movimiento_pendulo}
\end{equation}

Para pequeñas oscilaciones puede emplearse la aproximación

\begin{equation}
\sin\theta\approx\theta,
\label{eq:aproximacion_angular}
\end{equation}

con $\theta$ expresado en radianes. Bajo esta aproximación, la ecuación de movimiento toma la forma de un oscilador armónico simple:

\begin{equation}
\ddot{\theta}+\frac{g}{L}\theta=0.
\label{eq:oas_pendulo}
\end{equation}

La frecuencia angular natural del sistema queda definida por

\begin{equation}
\omega=\sqrt{\frac{g}{L}},
\label{eq:frecuencia_angular}
\end{equation}

y el período de oscilación viene dado por

\begin{equation}
T=2\pi\sqrt{\frac{L}{g}}.
\label{eq:periodo_pendulo}
\end{equation}

Esta expresión muestra que, dentro del modelo ideal y bajo la aproximación de pequeñas oscilaciones, el período depende de la longitud del péndulo y de la aceleración de la gravedad, pero no de la masa de la esfera. El período $T$ corresponde al tiempo requerido para completar un ciclo de oscilación. La amplitud angular inicial $\theta_0$ representa el desplazamiento angular máximo respecto a la posición de equilibrio al inicio del movimiento.

En un sistema real existen mecanismos disipativos, como la resistencia del aire y el rozamiento en el soporte. Una forma común de modelar estas pérdidas es mediante un término de amortiguamiento viscoso lineal proporcional a la velocidad. En este modelo, la fuerza disipativa se representa como

\begin{equation}
F_r=-bv,
\label{eq:fuerza_disipativa}
\end{equation}

donde $b$ es la constante de amortiguamiento. Para un régimen subamortiguado, la amplitud de las oscilaciones decrece progresivamente con el tiempo y puede modelarse mediante una envolvente exponencial de la forma

\begin{equation}
\theta_{\max}(t)=\theta_0e^{-\gamma t},
\label{eq:envolvente_amortiguada}
\end{equation}

donde

\begin{equation}
\gamma=\frac{b}{2m}.
\label{eq:constante_atenuacion}
\end{equation}

Este modelo permite describir la disminución gradual de la amplitud mientras el sistema continúa oscilando alrededor de la posición de equilibrio.
